(1) 换元 \(p=mc\sinh u\)，则 \(H=mc^2\cosh u\)，\(dp=mc\cosh u\,du\)。两种动量符号对称，所以
\[
Z_p=\int_\mathbb R e^{-\beta H}dp=2mcK_1(z).
\]
经典稀薄、交换量子效应可忽略时，
\[
\boxed{Z_N=\frac1{N!}\left(\frac{2mcL}{h}K_1(z)\right)^N},\qquad
f_T(p)=\frac{e^{-\beta H(p)}}{2mcK_1(z)}.
\]
若粒子可区分，应去掉Gibbs \(1/N!\)，而单粒子能量公式不变。热力学积分中包含静质量能，故低温配分仍具有因子 \(e^{-z}\)。

(2) \(K_1'=-K_0-K_1/z\)，设 \(R=K_0/K_1\)，求导
\[
\boxed{u=U/N=mc^2(R+z^{-1})},\qquad
\boxed{c_V=k_B[1+z^2(1-R^2)-zR]}.
\]
第二式来自 \(R'=-1+R^2+R/z\)，而
\(c_V=\operatorname{Var}(H)/(k_BT^2)>0\) 直接保证正性。
当 \(z\gg1\)，利用
\(K_\nu(z)\sim \sqrt{\pi/(2z)}e^{-z}[1+(4\nu^2-1)/(8z)+\cdots]\)，得
\(u=mc^2+k_BT/2+O((k_BT)^2/mc^2)\)，
\(c_V\to k_B/2\)。
当 \(z\ll1\)，\(K_1\sim 1/z,K_0\sim-\ln(z/2)-\gamma_E\)，得到 \(u\sim k_BT,c_V\to k_B\)，体现一维色散关系由 \(p^2/(2m)\) 变为 \(c|p|\)。

(3) 从 \(F=-k_BT\ln Z_N\) 对 \(L\) 求导有
\[
PL=Nk_BT,\qquad S=k_B(\ln Z_N+\beta U)
=k_B[\ln Z_N+N(zR+1)].
\]
固定 \(m,N\) 的准静态等熵过程满足
\[
\boxed{L K_1(z)e^{zR(z)+1}=\mathrm{const}}.
\]
大 \(z\) 时 \(K_1 e^{zR+1}\propto z^{-1/2}\)，所以
\(LT^{1/2}=\mathrm{const}\)；小 \(z\) 时 \(K_1\sim z^{-1}, zR\to0\)，故
\(LT=\mathrm{const}\)。
\(PL=Nk_BT\) 在两端形式不变，但热容和绝热幂律因能量色散而改变。

(4) 群速度为 \(v=dH/dp=pc^2/H\)。在 \(p>0\) 用 \(v\,dp=dH\)：
\[
\int_{\mathbb R}|v|e^{-\beta H}dp
=2\int_{mc^2}^\infty e^{-\beta E}dE=\frac{2e^{-z}}{\beta}.
\]
归一化给
\[
\boxed{\langle|v|\rangle=c\frac{e^{-z}}{zK_1(z)}}.
\]
低温 \(z\gg1\)，该式为 \(c\sqrt{2/(\pi z)}=
\sqrt{2k_BT/(\pi m)}\)，吻合一维Maxwell分布；
高温 \(z\ll1\) 时趋于 \(c\)。
对任意有限 \(p\)，\(|v(p)|<c\)，而正则分布对每个有限 \(p\) 有正的密度，因此任何有限温度下严格有 \(\langle|v|\rangle<c\)，没有超光速平均速度。
